% Generated by tools/build_new_dists_anthology_sections.py; do not edit by hand.

\section{Compact Spin and Pin groups}\label{proof:new-dist:spin_pin}

\resultbox{\textbf{Canonical testing artifacts.}\par\smallskip Exact backend: \path{new_dists/Spin_pin/verification.py}.\par Regression: \path{tests/test_spin_pin_new_dists.py}.\par Marimo: \path{marimo_notebooks/spin_pin.py}.}

\subsection*{Scope and outcome}

We prove the exact compact Molien--Weyl formulas for the spin representation
of $\operatorname{Spin}(2n+1)$ and the two half-spin representations of
$\operatorname{Spin}(2n)$.  Clifford algebra decompositions give the decisive
fourth multilinear coefficients $n+1$ and, respectively,
$n/2+1$ or $(n-1)/2$.  They grow with rank, so the unnormalized spinor
sigma-MGFs have no coefficientwise large-rank limit.  We also state the exact
two-component Pin formula, distinguish the $\operatorname{Pin}^{+}$ and
$\operatorname{Pin}^{-}$ lifts, and record the mixed spinor--vector and
adjoint consequences.  An accompanying marimo notebook independently
constructs Weyl characters and Clifford matrices.  All 76 exact and
matrix-level checks pass through rank five where applicable.

\subsection{The common compact coefficient formula}

Let $G$ be compact and let $W$ be a finite-dimensional unitary $G$-module
with representation $\rho$ and character $\chi_W$.  For a partition
$\tau=(\tau_1,\ldots,\tau_\ell)$, set
\[
 W_\tau=\bigotimes_{j=1}^{\ell}\Sym^{\tau_j}W,
 \qquad
 a_\tau(G,W)=\dim W_\tau^G.
\]
The sigma-MGF in the monomial basis is
\[
 \MGF_{G,W}(\mathbf t)=\sum_\tau a_\tau(G,W)m_\tau(\mathbf t).
\]

\begin{theorem}[compact generalized Molien formula]\label{nd:spin_pin:thm:molien}
For every compact $G$ and unitary $W$,
\begin{align}
 \MGF_{G,W}(\mathbf t)
 &=\int_G\prod_{i\ge1}\det(I-t_i\rho(g))^{-1}\,dg,\label{nd:spin_pin:eq:det-molien}\\
 &=\int_G\exp\!\left(\sum_{r\ge1}\frac{p_r(\mathbf t)}r
              \chi_W(g^r)\right)dg.\label{nd:spin_pin:eq:power-molien}
\end{align}
If $T$ is a maximal torus, $\Phi^+$ is a positive root system, and
$W_G$ is the Weyl group, then
\begin{equation}\label{nd:spin_pin:eq:weyl-master}
 \boxed{\MGF_{G,W}=\frac1{|W_G|}\CT\!\left[
 \prod_{\alpha\in\Phi^+}(1-z^\alpha)(1-z^{-\alpha})
 \exp\!\left(\sum_{r\ge1}\frac{p_r}{r}\chi_W(z^r)\right)
 \right].}
\end{equation}
\end{theorem}

\begin{proof}
Normalized Haar averaging is the Reynolds projection
$P_\tau=\int_G\rho_{W_\tau}(g)\,dg$.  It is an idempotent with image
$W_\tau^G$, so its trace is $a_\tau(G,W)$.  The identity
\[
 \sum_{d\ge0}\Tr(\Sym^d A)t^d=\det(I-tA)^{-1}
\]
and multiplication over the independent variables $t_i$ prove
\eqref{nd:spin_pin:eq:det-molien}.  Taking the formal logarithm of the determinant gives
\eqref{nd:spin_pin:eq:power-molien}.  Finally, the normalized Weyl integration formula
applied coefficient by coefficient gives \eqref{nd:spin_pin:eq:weyl-master}.
\end{proof}

For computation, let $\rho_\Phi$ be the Weyl vector and let $f$ be any
$W_G$-invariant Laurent polynomial.  The Weyl-denominator identity reduces
the two-sided denominator to the finite signed lookup
\begin{equation}\label{nd:spin_pin:eq:alternant-lookup}
 \frac1{|W_G|}\CT\!\left[
 \prod_{\alpha>0}(1-z^\alpha)(1-z^{-\alpha})f(z)\right]
 =\sum_{w\in W_G}\det(w)\,[z^{\rho_\Phi-w\rho_\Phi}]f(z).
\end{equation}
Equation \eqref{nd:spin_pin:eq:alternant-lookup} is the exact integer algorithm used by
the verifier.  Half-integral spin weights are stored in doubled orthonormal
coordinates, so no fractional exponent arithmetic is required.

\subsection{\texorpdfstring{Type $B_n$: $\operatorname{Spin}(2n+1)$ on its spin module}%
{Type Bn: Spin(2n+1) on its spin module}}

Let $V=V_{2n+1}$ be the vector representation and let $\Delta_n$ be the
$2^n$-dimensional complex spin representation.  In orthonormal torus
coordinates the weights are
\[
 \frac12(\epsilon_1e_1+\cdots+\epsilon_ne_n),
 \qquad \epsilon_j\in\{\pm1\},
\]
each once.  Hence
\begin{equation}\label{nd:spin_pin:eq:b-character}
 \chi_{\Delta_n}(z)=\prod_{j=1}^n(z_j^{1/2}+z_j^{-1/2}).
\end{equation}
Since $|W(B_n)|=2^nn!$, Theorem~\ref{nd:spin_pin:thm:molien} gives
\begin{equation}\label{nd:spin_pin:eq:b-mgf}
 \boxed{\MGF_{\operatorname{Spin}(2n+1),\Delta_n}
 =\frac1{2^nn!}\CT\!\left[\mathcal D_{B_n}(z)
 \exp\!\left(\sum_{r\ge1}\frac{p_r}{r}
 \prod_{j=1}^n(z_j^{r/2}+z_j^{-r/2})\right)\right].}
\end{equation}
Equivalently, if
\[
 H_k^{\Delta_n}(z)=[u^k]\exp\!\left(
 \sum_{r\ge1}\frac{u^r}{r}\chi_{\Delta_n}(z^r)\right),
\]
then every requested coefficient is the finite calculation
\begin{equation}\label{nd:spin_pin:eq:b-coeff}
 [m_\tau]\MGF
 =\frac1{2^nn!}\CT\!\left[
 \mathcal D_{B_n}(z)\prod_jH_{\tau_j}^{\Delta_n}(z)\right].
\end{equation}

The nontrivial element in the kernel of
$\operatorname{Spin}(2n+1)\to SO(2n+1)$ acts as $-I$ on $\Delta_n$.
It follows immediately that
\begin{equation}\label{nd:spin_pin:eq:b-parity}
 [m_\tau]\MGF=0\qquad\text{if $|\tau|$ is odd}.
\end{equation}
The spin module is self-dual.  The Clifford-periodic invariant bilinear form
is symmetric for $n\equiv0,3\pmod4$ and alternating for
$n\equiv1,2\pmod4$.  Therefore
\begin{equation}\label{nd:spin_pin:eq:b-quadratic}
 [m_{(1,1)}]\MGF=1,
 \qquad
 [m_{(2)}]\MGF=\begin{cases}
 1,&n\equiv0,3\pmod4,\\0,&n\equiv1,2\pmod4.
 \end{cases}
\end{equation}
One proof of the symmetry sign starts with the alternating form on the
$SU(2)\cong\operatorname{Spin}(3)$ module and tensors the Pauli-matrix
Clifford construction.  At each two-dimensional Clifford extension the
transpose sign changes in the periodic pattern $--++$.

\begin{theorem}[odd-spin fourth moment]\label{nd:spin_pin:thm:b-fourth}
For every $n\ge1$,
\begin{equation}\label{nd:spin_pin:eq:b-fourth}
 \boxed{\dim(\Delta_n^{\otimes4})^{\operatorname{Spin}(2n+1)}=n+1.}
\end{equation}
\end{theorem}

\begin{proof}
The even Clifford algebra acts faithfully on $\Delta_n$, and the symbol map
from Clifford products to exterior forms is Spin-equivariant.  Hodge duality
then gives the multiplicity-free decomposition
\begin{equation}\label{nd:spin_pin:eq:b-tensor-square}
 \Delta_n\otimes\Delta_n
 \cong\End(\Delta_n)
 \cong\bigoplus_{k=0}^{n}\Lambda^kV_{2n+1}.
\end{equation}
The summands are pairwise nonisomorphic and self-dual.  Since $\Delta_n$ is
self-dual,
\[
 \dim(\Delta_n^{\otimes4})^G
 =\dim\End_G(\Delta_n\otimes\Delta_n)=n+1.
\]
\end{proof}

\subsection{\texorpdfstring{Type $D_n$: half-spin modules of $\operatorname{Spin}(2n)$}%
{Type Dn: half-spin modules of Spin(2n)}}

Let $\Delta_n^+$ and $\Delta_n^-$ be the two half-spin modules, each of
dimension $2^{n-1}$.  Set
\[
 A_n(z)=\prod_{j=1}^n(z_j^{1/2}+z_j^{-1/2}),\qquad
 B_n(z)=\prod_{j=1}^n(z_j^{1/2}-z_j^{-1/2}).
\]
Separating spin weights by the parity of their minus signs gives
\begin{equation}\label{nd:spin_pin:eq:d-character}
 \chi_{\Delta_n^\pm}(z)=\frac12(A_n(z)\pm B_n(z)).
\end{equation}
Because $|W(D_n)|=2^{n-1}n!$,
\begin{equation}\label{nd:spin_pin:eq:d-mgf}
 \boxed{\MGF_{\operatorname{Spin}(2n),\Delta_n^\pm}
 =\frac1{2^{n-1}n!}\CT\!\left[\mathcal D_{D_n}(z)
 \exp\!\left(\sum_{r\ge1}\frac{p_r}{r}
 \chi_{\Delta_n^\pm}(z^r)\right)\right].}
\end{equation}

The half-spin central character has order two for even $n$ and order four
for odd $n$.  Thus
\begin{equation}\label{nd:spin_pin:eq:d-selection}
 [m_\tau]\MGF=0\quad
 \begin{cases}
 \text{if $|\tau|$ is odd},&n\text{ even},\\
 \text{unless $4\mid|\tau|$},&n\text{ odd}.
 \end{cases}
\end{equation}
Moreover,
\[
 (\Delta_n^+)^*\cong
 \begin{cases}\Delta_n^+,&n\text{ even},\\
 \Delta_n^-,&n\text{ odd},\end{cases}
\]
and, in the self-dual case, the form is symmetric for $n\equiv0\pmod4$
and alternating for $n\equiv2\pmod4$.  Consequently
\begin{equation}\label{nd:spin_pin:eq:d-quadratic}
 [m_{(1,1)}]\MGF=\mathbf1_{2\mid n},
 \qquad
 [m_{(2)}]\MGF=\mathbf1_{4\mid n}.
\end{equation}

\begin{theorem}[half-spin fourth moments]\label{nd:spin_pin:thm:d-fourth}
For $n\ge2$,
\begin{equation}\label{nd:spin_pin:eq:d-fourth}
 \boxed{\dim((\Delta_n^+)^{\otimes4})^{\operatorname{Spin}(2n)}
 =\begin{cases}
 n/2+1,&n\text{ even},\\[2pt]
 (n-1)/2,&n\text{ odd}.
 \end{cases}}
\end{equation}
The balanced count is
\begin{equation}\label{nd:spin_pin:eq:d-balanced}
 \boxed{\dim\left((\Delta_n^+\otimes\Delta_n^{+*})^{\otimes2}
 \right)^{\operatorname{Spin}(2n)}=\left\lfloor\frac n2\right\rfloor+1.}
\end{equation}
\end{theorem}

\begin{proof}
The Clifford symbol map and the two Hodge eigenspaces in middle degree give
\begin{equation}\label{nd:spin_pin:eq:d-tensor-square}
 \Delta_n^+\otimes\Delta_n^+
 \cong
 \bigoplus_{\substack{0\le k<n\\k\equiv n\ ({\rm mod}\ 2)}}
 \Lambda^kV_{2n}\ \oplus\ \Lambda^n_{\varepsilon}V_{2n},
\end{equation}
where $\Lambda^n_{\varepsilon}$ is one of the two middle-degree irreducibles.
The tensor square of $\Delta_n^-$ has the same lower-degree summands and the
opposite middle constituent.

If $n$ is even, $\Delta_n^+$ is self-dual and
\eqref{nd:spin_pin:eq:d-tensor-square} has $n/2+1$ pairwise nonisomorphic summands.
The invariant fourth power is its endomorphism algebra, proving the first
line of \eqref{nd:spin_pin:eq:d-fourth}.  If $n$ is odd, invariants are
\[
 \Hom_G(\Delta_n^-\otimes\Delta_n^-,
         \Delta_n^+\otimes\Delta_n^+).
\]
Only the $(n-1)/2$ lower exterior powers occur on both sides; the middle
constituents are opposite.  This proves the second line.

Finally, $\Delta_n^+\otimes\Delta_n^{+*}=\End(\Delta_n^+)$ is
multiplicity-free.  Its exterior-form constituents are the even degrees below
$n$, together with one middle summand when $n$ is even.  Their number is
$\lfloor n/2\rfloor+1$, proving \eqref{nd:spin_pin:eq:d-balanced}.
\end{proof}

The special case $n=3$ is a useful boundary check:
$\operatorname{Spin}(6)\cong SU(4)$ and $\Delta_3^+$ is the defining
four-dimensional representation.  Its first one-copy invariant is the
degree-four determinant, so the fourth multilinear coefficient is one, while
both quadratic coefficients vanish.

\subsection{Asymptotic obstruction}

\begin{corollary}\label{nd:spin_pin:cor:no-limit}
Neither the unnormalized odd-spin sequence
$\MGF_{\operatorname{Spin}(2n+1),\Delta_n}$ nor either half-spin sequence has
a finite coefficientwise limit in the completed symmetric-function ring.
\end{corollary}

\begin{proof}
The fixed coefficient $[m_{(1,1,1,1)}]$ equals $n+1$ in type $B_n$ and is
asymptotic to $n/2$ in type $D_n$.  It diverges along every sequence with
$n\to\infty$.  Coefficientwise convergence would require this scalar
coefficient to converge, which is impossible.
\end{proof}

The conclusion is not contradicted by $\dim\Delta_n=2^n$.  Increasing rank
opens a new exterior-power channel in \eqref{nd:spin_pin:eq:b-tensor-square} or every two
ranks in \eqref{nd:spin_pin:eq:d-tensor-square}.  A normalized, logarithmic, or cumulant
renormalization may still have a limit, but the raw sigma-MGF does not.

\subsection{Pin groups: the disconnected component formula}

Fix a concrete complex pinor representation $\Pi_N^\varepsilon$ of
$\operatorname{Pin}^\varepsilon(N)$, $\varepsilon\in\{+,-\}$, and choose
$u_\varepsilon$ in the odd component.  Since
\[
 \operatorname{Pin}^\varepsilon(N)
 =\operatorname{Spin}(N)\sqcup
 u_\varepsilon\operatorname{Spin}(N),
\]
normalized Haar measure assigns mass $1/2$ to each component.  Therefore
\begin{align}\label{nd:spin_pin:eq:pin-components}
 \MGF_{\operatorname{Pin}^\varepsilon(N),\Pi_N^\varepsilon}
 ={}&\frac12\int_{\operatorname{Spin}(N)}
 \exp\!\left(\sum_{r\ge1}\frac{p_r}{r}
 \chi_{\Pi_N^\varepsilon}(g^r)\right)dg\\
 &+\frac12\int_{\operatorname{Spin}(N)}
 \exp\!\left(\sum_{r\ge1}\frac{p_r}{r}
 \chi_{\Pi_N^\varepsilon}((u_\varepsilon g)^r)\right)dg.\notag
\end{align}
This is an exact identity, but the second integrand depends on the chosen Pin
extension, not only on the restricted Spin character.

For $N=2n$, the restriction is $\Delta_n^+\oplus\Delta_n^-$.  An odd
Clifford generator anticommutes with chirality and is therefore block
off-diagonal.  Every odd power remains block off-diagonal, so
\begin{equation}\label{nd:spin_pin:eq:pin-odd-trace}
 \chi_{\Pi_{2n}^\varepsilon}(x^{2r+1})=0
 \qquad(x\notin\operatorname{Spin}(2n)).
\end{equation}
Only the $p_{2r}$ terms remain in the odd-component exponential.  The two Pin
groups must not be identified: in the concrete Clifford models used here,
\begin{equation}\label{nd:spin_pin:eq:pin-square}
 u_+=\gamma_1,\quad u_- = i\gamma_1,
 \qquad u_+^2=I,\quad u_-^2=-I.
\end{equation}
Thus their even power characters can differ.

\subsection{Mixed spinor--vector and adjoint modules}

Character calculus immediately extends the exact constant term.  For
$W_{a,b}=\Delta^{\otimes a}\otimes V^{\otimes b}$,
\begin{equation}\label{nd:spin_pin:eq:mixed-character}
 \chi_{W_{a,b}}(z)=\chi_\Delta(z)^a\chi_V(z)^b,
\end{equation}
where
\[
 \chi_{V_{2n+1}}=1+\sum_j(z_j+z_j^{-1}),\qquad
 \chi_{V_{2n}}=\sum_j(z_j+z_j^{-1}).
\]
Substitution of \eqref{nd:spin_pin:eq:mixed-character} into \eqref{nd:spin_pin:eq:weyl-master} is the
exact sigma-MGF at every finite rank.  For a direct sum $\Delta\oplus V$, the
character product is replaced by $\chi_\Delta+\chi_V$.

The adjoint representation factors through $SO(N)$ and is
$\mathfrak{so}_N\cong\Lambda^2V_N$.  Its torus character is
\begin{equation}\label{nd:spin_pin:eq:adjoint-character}
 \chi_{\rm Ad}(z)=\operatorname{rank}(G)+\sum_{\alpha\in\Phi}z^\alpha,
\end{equation}
so \eqref{nd:spin_pin:eq:weyl-master} again gives an exact finite-rank formula.  The
exterior-square identity gives the directly testable trace formula
\begin{equation}\label{nd:spin_pin:eq:adjoint-trace}
 \boxed{\Tr(\operatorname{Ad}(g)^r)=\frac12\left(
 \Tr(V(g)^r)^2-\Tr(V(g)^{2r})\right).}
\end{equation}
Unlike the genuine spinors, $\Lambda^2V_N$ has a coefficientwise stable
orthogonal regime.  For fixed $\tau$ and the safe range $N>2|\tau|$,
\begin{equation}\label{nd:spin_pin:eq:adjoint-stable}
 [m_\tau]\MGF_{{\rm Ad},\infty}
 =\left\langle\ExpS(h_2),h_\tau[e_2]\right\rangle.
\end{equation}
The strict inequality is a sufficient stable range, not a claim of
sharpness in every degree.

\subsection{Finite restrictions and spin covers}

If $H\le\operatorname{Spin}(N)$ is finite and $W$ is a restricted spinor
module, Reynolds averaging gives
\begin{equation}\label{nd:spin_pin:eq:finite-restriction}
 \MGF_{H,W}=\frac1{|H|}\sum_{h\in H}
 \exp\!\left(\sum_{r\ge1}\frac{p_r}{r}\chi_W(h^r)\right)
 =\sum_{[h]}\frac1{|C_H(h)|}
 \exp\!\left(\sum_{r\ge1}\frac{p_r}{r}\chi_W(h^r)\right).
\end{equation}
The same formula applies to the inverse image $\widetilde W$ of a finite real
reflection group in Pin and to a Schur double cover $\widetilde S_n$ with a
chosen genuine spin character.  The cover and power maps are essential:
ordinary conjugacy classes can split upstairs.  No representation-independent
large-$n$ law exists for $\widetilde S_n$ until a sequence of strict-partition
labels is specified.

\subsection{Independent computational verification}

\subsubsection{Exact Weyl route}

For type $B_n$ the verifier enumerates all signed permutations; for type
$D_n$ it retains those with an even number of sign changes.  In doubled
coordinates,
\[
 2\rho_{B_n}=(2n-1,2n-3,\ldots,1),\qquad
 2\rho_{D_n}=(2n-2,2n-4,\ldots,0).
\]
For each $\tau$, dynamic programming expands
$\prod_jh_{\tau_j}$ from the actual spin weights.  The finite lookup
\eqref{nd:spin_pin:eq:alternant-lookup} then returns the invariant multiplicity.  The
decompositions in Theorems~\ref{nd:spin_pin:thm:b-fourth} and \ref{nd:spin_pin:thm:d-fourth} are not
used by this route.

\begin{center}
\begin{tabular}{@{}c@{\quad}r@{\quad}rrrr@{\quad}c@{}}
\toprule
Group&$\dim W$&$[m_{(2)}]$&$[m_{(1,1)}]$&$[m_{(1^4)}]$&formula&result\\
\midrule
$\operatorname{Spin}(3)$&2&0&1&2&$n+1$&\textcolor{teal}{\textbf{PASS}}\\
$\operatorname{Spin}(5)$&4&0&1&3&$n+1$&\textcolor{teal}{\textbf{PASS}}\\
$\operatorname{Spin}(7)$&8&1&1&4&$n+1$&\textcolor{teal}{\textbf{PASS}}\\
$\operatorname{Spin}(9)$&16&1&1&5&$n+1$&\textcolor{teal}{\textbf{PASS}}\\
$\operatorname{Spin}(11)$&32&0&1&6&$n+1$&\textcolor{teal}{\textbf{PASS}}\\
\bottomrule
\end{tabular}
\end{center}

\begin{center}
\begin{tabular}{@{}c@{\quad}r@{\quad}rrrr@{\quad}c@{}}
\toprule
Group&$\dim\Delta^+$&$[m_{(2)}]$&$[m_{(1,1)}]$&$[m_{(1^4)}]$&balanced fourth&result\\
\midrule
$\operatorname{Spin}(4)$&2&0&1&2&2&\textcolor{teal}{\textbf{PASS}}\\
$\operatorname{Spin}(6)$&4&0&0&1&2&\textcolor{teal}{\textbf{PASS}}\\
$\operatorname{Spin}(8)$&8&1&1&3&3&\textcolor{teal}{\textbf{PASS}}\\
$\operatorname{Spin}(10)$&16&0&0&2&3&\textcolor{teal}{\textbf{PASS}}\\
\bottomrule
\end{tabular}
\end{center}

All odd-total-degree type-$B$ checks vanish.  The type-$D$ suite separately
checks the order-two/order-four selection rule.  In total, 45 exact spin,
half-spin, and balanced coefficient comparisons pass.

\subsubsection{Direct Clifford-matrix route}

Let $X,Y,Z$ be the Pauli matrices.  The implementation uses the
Jordan--Wigner gamma matrices
\[
 \gamma_{2j+1}=Z^{\otimes j}\otimes X\otimes I^{\otimes(n-j-1)},\qquad
 \gamma_{2j+2}=Z^{\otimes j}\otimes Y\otimes I^{\otimes(n-j-1)},
\]
with $Z^{\otimes n}$ appended in odd dimension.  They satisfy
\[
 \gamma_a\gamma_b+\gamma_b\gamma_a=2\delta_{ab}I.
\]
Concrete represented Spin elements are products
\begin{equation}\label{nd:spin_pin:eq:matrix-spin}
 \rho(g)=\prod_{a<b}\exp\!\left(
 \frac{\theta_{ab}}2\gamma_a\gamma_b\right).
\end{equation}
Because $(\gamma_a\gamma_b)^2=-I$, each exponential in
\eqref{nd:spin_pin:eq:matrix-spin} is evaluated exactly as a sine--cosine matrix
polynomial.

On a tensor square, the code constructs the positive quadratic Casimir
\[
 C=-\sum_{a<b}\left(J_{ab}\otimes I+I\otimes J_{ab}\right)^2,
 \qquad J_{ab}=\tfrac12\gamma_a\gamma_b.
\]
Its eigenspace dimensions provide an independent matrix signature of the
irreducible summands.

\begin{center}\small
\begin{tabularx}{\linewidth}{@{}lrrXc@{}}
\toprule
Module&matrix dimension&Casimir blocks&observed block dimensions&result\\
\midrule
$\operatorname{Spin}(3)$ spin&2&2&$1,3$&\textcolor{teal}{\textbf{PASS}}\\
$\operatorname{Spin}(5)$ spin&4&3&$1,5,10$&\textcolor{teal}{\textbf{PASS}}\\
$\operatorname{Spin}(7)$ spin&8&4&$1,7,21,35$&\textcolor{teal}{\textbf{PASS}}\\
$\operatorname{Spin}(9)$ spin&16&5&$1,9,36,84,126$&\textcolor{teal}{\textbf{PASS}}\\
$\operatorname{Spin}(4)$ half-spin&2&2&$1,3$&\textcolor{teal}{\textbf{PASS}}\\
$\operatorname{Spin}(8)$ half-spin&8&3&$1,28,35$&\textcolor{teal}{\textbf{PASS}}\\
\bottomrule
\end{tabularx}
\end{center}

The type-$B$ dimensions are exactly
$\binom{2n+1}{0},\ldots,\binom{2n+1}{n}$, as predicted by
\eqref{nd:spin_pin:eq:b-tensor-square}.  Every Clifford anticommutator residual is zero
in floating arithmetic for these Pauli matrices.

\subsubsection{Pin, mixed, and adjoint checks}

For $N=2,4,6$, the verifier constructs $u_+=\gamma_1$ and
$u_-=i\gamma_1$, a nontrivial Spin torus element, and their product in the
odd component.  All six cases have the correct square, anticommute with
chirality, and have zero traces in powers $1,3,5$.  The largest unitarity
residual is $5.7\times10^{-16}$.

Mixed spinor--vector Kronecker matrices in types $B_3$ and $D_3$ agree with
the product-weight construction and \eqref{nd:spin_pin:eq:mixed-character}.  Explicit
adjoint torus matrices verify \eqref{nd:spin_pin:eq:adjoint-trace} for $r=1,2,3,4$ with
maximum residual $7.2\times10^{-15}$.  Exact Weyl integration in types
$B_2,B_3,D_4$ also recovers one invariant metric and one invariant
multilinear structure tensor.  These are 25 additional checks.  Together
with the six Casimir rows, the complete report contains 76 passing checks.

\subsection{Reproduction and scope}

From the repository root, run
\begin{verbatim}
.venv/bin/python new_dists/verification.py
.venv/bin/python -m pytest -q tests/test_spin_pin_new_dists.py
.venv/bin/marimo edit marimo_notebooks/spin_pin.py
\end{verbatim}
The first command prints the six verification subtotals and fails with a
nonzero exit status if any comparison fails.  The notebook exposes the same
suite with a rank selector and group-separated tables.

The exact coefficient computations establish the tested low-degree cases;
the all-rank claims are proved by Reynolds/Weyl integration and the Clifford
decompositions.  The matrix computations are an independent finite-rank
audit, not a substitute for those proofs.  For Pin, the document proves the
component reduction and the even-dimensional odd-trace simplification.  A
fully expanded Pin sigma-MGF still requires the power characters of the
chosen $\operatorname{Pin}^{\pm}$ extension.

\begin{thebibliography}{9}
\bibitem{nd:spin_pin:GKRS}
B.~Gross, B.~Kostant, P.~Ramond, and S.~Sternberg,
\emph{The Weyl character formula, the half-spin representations, and equal
rank subgroups}, Proc. Natl. Acad. Sci. USA 95 (1998),
\href{https://arxiv.org/abs/math/9808133}{arXiv:math/9808133}.

\bibitem{nd:spin_pin:Koike}
K.~Koike,
\emph{Spin representations and centralizer algebras for the spinor groups},
\href{https://arxiv.org/abs/math/0502397}{arXiv:math/0502397}.

\bibitem{nd:spin_pin:Wenzl}
H.~Wenzl,
\emph{On centralizer algebras for spin representations},
\href{https://arxiv.org/abs/1107.4183}{arXiv:1107.4183}.

\bibitem{nd:spin_pin:LM}
H.~B.~Lawson, Jr. and M.-L.~Michelsohn,
\emph{Spin Geometry}, Princeton Mathematical Series 38,
Princeton University Press, 1989.
\end{thebibliography}
